\section{Introduction}
\label{sec:introduction}

A finalized Solana mainnet transaction verified a 75,358-byte transparent
proof of a private spend, recorded the public nullifier, and replaced the pool
root in one atomic state transition.  The verifier used 1,334,452 of the
transaction's 1,356,912 compute units.  Its commitments are hash based and
require no trusted setup.  To our knowledge, following a public-evidence
search completed on 24 August 2026, this is the first publicly evidenced
Solana mainnet transaction to combine direct verification of a transparent,
computationally hiding private-spend proof with the corresponding nullifier
and pool-state update in the same transaction.  The comparison covers publicly
indexed papers, repositories, deployments, and patents available by the cutoff
date.  The proof account was uploaded and sealed beforehand; the reported
transaction performed proof verification and the atomic spend-state update.
\Cref{sec:related-work} gives the closest public systems and exact comparison.

This paper presents that result together with a formal argument reaching from
the deployed verifier to its cryptographic meaning.  The central theorem
begins with an arbitrary successful call to the deployed Rust proof checker,
after functional translation into Lean~4.  From that call alone it derives
the parsed statement, every byte absorbed by Fiat--Shamir, all six work
checks, the eighteen distinct queries, five authenticated opening sections,
four low-degree folds, the final polynomial, all algebraic claims and relation
messages, and both final accumulator equalities.  These values belong to one
execution.  The theorem therefore rules out a real composition error: proving
the transcript for one accepted run and the openings or algebra for another
does not establish anything about either run.

The span of this theorem is the principal formal contribution.  Formal
cryptographic developments commonly prove an abstract protocol or a
collection of executable components.  Here the accepted-execution theorem
crosses the parser, transcript, authenticated data, low-degree test, algebraic
relation, and terminal arithmetic of a deployed blockchain verifier.  It is
paired with reproducible evidence connecting the reviewed Rust source to the
program bytes and those bytes to the finalized transaction.  The result is a
source-to-security argument with a concrete chain execution at its far end.

\subsection{The mathematical argument}

The proof system establishes a one-input, one-output spend.  Its public
statement contains the old and new pool roots, a one-time nullifier, the
output-note commitment, asset and fee data, and the deployment domain.  The
witness opens the input and output notes, proves ownership and integer value
conservation, and supplies a depth-twenty path that replaces the input leaf by
the output commitment.  We formalize this relation directly and prove that a
valid authenticated algebraic trace constructs such a witness.

The low-degree argument is developed for the exact verifier rather than an
asymptotic profile.  Lean proves the order of the chosen generator in
\(\mathbb F_{2^{31}-1}\), distinctness and ordering of the circle and line
domains, and the evaluation identities of all four encoders.  The initial
decoder list has size at most 240 under the published circle-decoding result.
Outside the counted challenge events, one member of that list follows a
single coefficient chain
\[
  1024 \longrightarrow 256 \longrightarrow 64
  \longrightarrow 16 \longrightarrow 4
\]
through the four radix-four folds and reaches the polynomial published by the
prover.  A separate causal argument fixes each exceptional challenge set from
the transcript prefix available before that challenge is sampled.

The query calculation also follows the implemented sampler exactly.  It
retains the first eighteen distinct positions found in at most sixty-four
draws.  Conditional on success, the output is uniform over ordered injections.
Thus, for a bad set of size \(A\) in a domain of size \(N\), the exact miss
probability is
\[
  \frac{(A)_{18}}{(N)_{18}}
  =\prod_{j=0}^{17}\frac{A-j}{N-j}.
\]
Here \((m)_k=m(m-1)\cdots(m-k+1)\) is the falling factorial.
The same expression bounds the unconditioned event because sampler exhaustion
rejects.  Finally, the relation proof shows algebraically that folding the
candidate and folding the verifier's dual weights preserve their dot product.
This carries the extracted polynomial into the spend constraints rather than
leaving low degree and statement validity as separate claims.

The formal results have the following logical shape.  For one translated
execution \(\rho\),
\begin{equation}
  \mathsf{Accept}(\rho)
  \Longrightarrow \mathsf{Evidence}(\rho),
  \label{eq:intro-accepted-evidence}
\end{equation}
where \(\mathsf{Evidence}(\rho)\) contains the full transcript,
authentication, folding, and relation data from \(\rho\).  In the security
experiment, Lean separately proves the pointwise inclusion
\begin{equation}
  A_{\mathrm{false}}
  \subseteq C\cup\bigcup_{i=1}^{18}E_i,
  \label{eq:intro-failure-coverage}
\end{equation}
where \(C\) contains the explicit protocol events and the \(E_i\) are named
primitive, credential, implementation, and runtime events.  Two displayed
interfaces connect these endpoints: fieldwise agreement between the live Rust
statement and the abstract statement, and a causal embedding of the
deterministic accepted-call classification into the probability experiment.

\subsection{End-to-end formal evidence}

We use \emph{end-to-end} for a theorem and artifact chain in which every
change of representation is exposed to review.  The evidence has five parts:

\begin{enumerate}[leftmargin=*]
  \item \textbf{Mathematical specification.}  Lean definitions state the
        private-spend relation, concrete domains and encoders, transcript,
        authenticated openings, low-degree test, algebraic relation, security
        experiment, and state transition.
  \item \textbf{Implementation proof.}  The deployed Rust proof-checking path
        is functionally translated into Lean.  Refinement theorems connect its
        control flow and returned values to the mathematical definitions.
  \item \textbf{Accepted-execution composition.}  One universal theorem
        derives every acceptance-critical fact, including both terminal dot
        products, from one successful translated call.
  \item \textbf{Program identity.}  Pinned source, dependencies, compiler
        tools, and a deterministic build reproduce the deployed program byte
        for byte.
  \item \textbf{Mainnet evidence.}  The archived transaction binds the
        program, proof, statement, compute usage, nullifier, and old and new
        pool states.
\end{enumerate}

\Cref{tab:evidence-pipeline} records the evidence supplied for each layer.
The manuscript uses mathematical names; the artifact guide maps them to exact
Lean declarations, files, commands, revisions, and hashes.

\begin{table}[t]
\centering
\small
\caption{Evidence supporting the end-to-end result.}
\label{tab:evidence-pipeline}
\begin{tabularx}{\textwidth}{@{}R{0.20\textwidth}YY@{}}
\toprule
Artifact & Established fact & Reproduction evidence \\
\midrule
Lean development & Spend mathematics, functional correctness, failure
decomposition, and the accepted-execution theorem & Clean replay of the
complete 331-module dependency closure \\
Rust-to-Lean translation & Control flow and values of the successful
proof-checking path & Pinned translator revision, generated Lean, and source
manifests \\
Build record & Identity between recorded source inputs and deployed program
bytes & Deterministic build command, toolchain record, and byte hashes \\
Mainnet record & Direct proof verification and atomic nullifier and pool
update & Signed transaction, account snapshots, proof bytes, and verification
script \\
\bottomrule
\end{tabularx}
\end{table}

The accepted-execution theorem is universal over successful calls, rather
than a replay theorem for the archived proof.  It reconstructs the public
input digest and transcript, verifies six sequential work conditions, obtains
eighteen distinct query positions, authenticates five opening sections,
checks four low-degree rounds and the final polynomial, decodes seventy-six
algebraic claims and fifty-eight relation fields, and derives the terminal
equalities of twelve-component and four-component accumulators.  Its
conclusion depends on the derived mathematical evidence, including the values
consumed before the last helper returns.

\subsection{Security statement}

All false-acceptance events live in one finite probability space.  The
coverage theorem in \cref{eq:intro-failure-coverage} is pointwise, so its
probability bound uses a union bound without an independence hypothesis.
For the evaluated parameters, Lean checks that the work-normalized protocol
contribution satisfies
\[
  B_{\mathrm{protocol}}\leq 0.7\cdot2^{-100}.
\]
Published circle-decoding and Fiat--Shamir theorems enter through explicit
interfaces whose concrete parameter hypotheses are checked in Lean.  If the
remaining primitive and credential events contribute at most
\(0.3\cdot2^{-100}\), and the stated deterministic correspondences hold, the
combined accepted-false probability is at most \(2^{-100}\).  A separate raw
resource bound gives no security credit for grinding work that has already
been completed; its dominant explicit term is approximately \(2^{-70}\).

\subsection{Contributions}

The paper makes four contributions:
\begin{itemize}[leftmargin=*]
  \item The first publicly evidenced Solana mainnet result, to our knowledge
        at the stated cutoff, that directly verifies a transparent,
        computationally hiding private-spend proof and atomically records its
        nullifier and new pool state within the transaction compute limit.
  \item An end-to-end Lean theorem for every successful translated execution
        of the deployed proof-checking path, covering the full chain from
        proof bytes and transcript operations through authenticated openings,
        low-degree folds, the algebraic relation, and both final accumulators.
  \item A complete mathematical treatment of the evaluated verifier:
        private-spend extraction, exact field and encoder results, coherent
        four-fold extraction, exact capped-query sampling, fold duality,
        theft and replay reductions, and the finite security experiment with
        checked concrete arithmetic.
  \item A reproducible evidence package connecting the theorem to pinned
        translation inputs, byte-identical program reproduction, proof and
        statement bytes, and a finalized mainnet state transition.
\end{itemize}

\subsection{Organization}

\Cref{sec:background} defines the spend relation and proof parameters.
\Cref{sec:lean-model} develops the checked mathematics, followed by the proof
method and implementation theorem in
\cref{sec:proof-design,sec:implementation}.  \Cref{sec:soundness} gives the
single security experiment and numerical bounds.
\Cref{sec:theft-state,sec:privacy} cover theft, state, and hiding, and
\cref{sec:evaluation} presents the formal replay, program reproduction, and
mainnet evidence.  Assumptions and related work follow; the appendices give
proof details and the independent reproduction checklist.
